\documentclass{ximera}

\begin{document}

    \title{Lecture 5}
    
    \begin{abstract}  
    Outline: \\
    - divergence \\
    - algebraic and order limit theorems
\end{abstract}

\maketitle

\begin{definition}
    A sequence that does not converge is said to diverge.
\end{definition}

\begin{definition}
    A sequence $\{a_n\}$ is bounded if there exists $ M \in \mathbb{R}_{>0}$ s.t. $|a_n| \leq M$ for all $n \in \mathbb{N}$.
\end{definition}

\begin{exercise}
    Show that, if a sequence has a limit, then the limit is unique.
\end{exercise}

\begin{theorem}
    A sequence that converges is bounded.
\end{theorem}
\begin{proof}
    Let $\{a_n\}$ be our sequence that converges to $a$.  Then, for all $\epsilon > 0$, there exists an $N \in \mathbb{N}$ such that $|a_n -a| < \epsilon$ for all $n \geq N$. Then, $a_n -a < \epsilon$ implies $a_n < a + \epsilon \leq |a| + \epsilon $, which implies that $|a_n| < |a|+ \epsilon$. Consider the set $\{a_1, \;\ldots, a_{N-1}\}$.  Then, for all $n \in \mathbb{N}$
    \begin{align*}
        |a_n| \leq \max \left\{a_1, \;\ldots, a_{N-1}, |a| +\epsilon \right\}.
    \end{align*}
    Therefore, the sequence is bounded.
\end{proof}

\textcolor{red}{(Have them conjecture about the converse and provide an example)}

\begin{exercise}
    Suppose $\lim a_n = a$, $\lim b_n = b$.
    Prove 
    \begin{enumerate}
        \item $\lim(ca_n) = ca$ $\; \forall \;c \in \mathbb{R}$
        \item $\lim (a_n + b_n) = a+b$
        \item (Bonus) $\lim(a_nb_n) = ab$
    \end{enumerate}
\end{exercise}

\begin{theorem}{(Algebraic Limit Theorem)}
Suppose $\lim a_n = a$, $\lim b_n = b$.
    \begin{itemize}
        \item $\lim(ca_n) = ca$ $\; \forall \;c \in \mathbb{R}$
        \item $\lim (a_n + b_n) = a+b$
        \item $\lim(a_nb_n) = ab$
        \item $\lim(a_n/b_n) = a/b$ if $b \neq 0$
    \end{itemize}
\end{theorem}

\begin{exercise}
    Prove that $\lim \frac{2n}{n+1} = 2$ two ways: using the definition of convergence only, and then using the Algebraic Limit Theorem. 
\end{exercise}

\begin{proof}
    Here is the proof using the definition only.  Let $\epsilon > 0$
    Consider 
    \begin{align}
        \left|\frac{2n}{n+1} -2\right| &= \left|\frac{2n}{n+1} - \frac{2(n+1)}{n+1} \right|
        \\
        &= \left|\frac{-2}{n+1}\right|
        \\
        &= \frac{2}{n+1}.
    \end{align}
    Then, notice that by the Archimedean property(b), there exists $N \in \mathbb{N}$ such that $\frac{1}{N+1} < \frac{1}N < \epsilon/2$. Then, $\frac{2}{N+1} < \epsilon$, and then by the properties of the natural numbers, $\frac{2}{n+1} < \epsilon$ for all $n \geq N$.  
\end{proof}

\newpage
\begin{theorem}{(Order Limit Theorem)}
    Let $\lim a_n = a$, $\lim b_n = b$.
    \begin{itemize}
        \item If $a_n \geq 0$ for all $n \in \mathbb{N}$, then $a \geq 0$.
        \item If $a_n \leq b_n$ for all $n \in \mathbb{N}$, then $a \leq b$.
        \item If there exists $c \in \mathbb{R}$ such that $c \leq b_n$, then $c \leq b$.  Similarly, if $d \in \mathbb{R}$ such that $a_n \leq d$, then $a \leq d$.
    \end{itemize}
\end{theorem}

\begin{definition}
    A sequence $\{a_n\}$ is increasing if $a_n \leq a_{n+1}$ for all $n \in \mathbb{N}$, and decreasing if $a_n \geq a_{n+1}$ for all $n \in \mathbb{N}$.  A sequence is called monotonic if it is either increasing or decreasing.
\end{definition}

\textcolor{red}{Have them think for a minute about what might happen if a sequence is monotone and bounded.}

\begin{theorem}{(Monotone Convergence Theorem)}
    If a sequence is monotone and bounded, then it converges.
\end{theorem}
Prove next time.

\end{document}